% ==========================================================================
% Взвешенная задача Горейнова--Тыртышникова--Замарашкина.
% Что доказано, что осталось, и как идти дальше.
% Сборка:  pdflatex brief.tex  (дважды)
% ==========================================================================
\documentclass[11pt,a4paper]{article}

\usepackage[T2A]{fontenc}
\usepackage[utf8]{inputenc}
\usepackage[russian]{babel}
\usepackage{lmodern}
\usepackage[left=25mm,right=25mm,top=24mm,bottom=26mm]{geometry}
\setlength{\emergencystretch}{3em}
\usepackage{amsmath,amssymb,amsthm}
\usepackage{microtype}
\usepackage{array}
\usepackage{longtable}
\usepackage{xcolor}
\usepackage[colorlinks=true,linkcolor=blue!50!black,urlcolor=blue!50!black,
            citecolor=blue!50!black]{hyperref}

\theoremstyle{plain}
\newtheorem{theorem}{Теорема}[section]
\newtheorem{proposition}[theorem]{Предложение}
\newtheorem{lemma}[theorem]{Лемма}
\newtheorem{corollary}[theorem]{Следствие}
\newtheorem{conjecture}[theorem]{Гипотеза}
\theoremstyle{definition}
\newtheorem{example}[theorem]{Пример}
\newtheorem{task}[theorem]{Задача}
\theoremstyle{remark}
\newtheorem{remark}[theorem]{Замечание}

\DeclareMathOperator{\tr}{tr}
\DeclareMathOperator{\adj}{adj}
\DeclareMathOperator{\rank}{rank}
\newcommand{\RR}{\mathbb{R}}
\newcommand{\CC}{\mathbb{C}}
\newcommand{\FF}{\mathbb{F}}
\newcommand{\Id}{\mathrm{I}}
\newcommand{\GTZ}{\mathrm{GTZ}}
\newcommand{\Sk}{\mathcal{S}}
\newcommand{\code}[1]{\texttt{\small #1}}
\newcommand{\T}{^{\!\top}}

\title{Взвешенная задача Горейнова--Тыртышникова--Замарашкина:\\
что доказано, что осталось и как идти дальше}
\author{}
\date{}

\begin{document}
\maketitle

\begin{abstract}
\noindent
Даны $m$ векторов $g_c\in\RR^k$ и положительные веса $t_c$, $\sum_c t_c=1$,
такие что $\sum_c t_c\,g_cg_c\T=\Id$. Спрашивается, найдутся ли среди них
$k$ векторов с $\sum_{c\in C}g_cg_c\T\succeq\Id$. Это взвешенная форма вопроса
о подматрице $B$ матрицы с ортонормированными столбцами, у которой
$\|B^{-1}\|_2\le\sqrt m$. Доказано: ответ «да» при $m\le k+2$ (любой ранг) и
при $k\le2$ (любой размер); над $\CC$ ответ «да» ровно при $m\le k+1$.
Остаётся одна клетка, $m=6$, $k=3$, и она решает весь ранг~3. Клетка сведена
к одному утверждению о «граничных» системах: у каждой есть пара
коллинеарных векторов. Оно, в свою очередь, разобрано на короткий список
конкретных полиномиальных задач. Ниже: словарь, известное, разбор
клетки $(6,3)$ с точным списком остатка, и дорога от $(6,3)$ к рангу~3
и к произвольному рангу. Утверждения с именами в \code{моноширинном}
шрифте формализованы в Lean~4 и проверены ядром; утверждения без имени
доказаны в тексте.
\end{abstract}

\tableofcontents

% ==========================================================================
\section{Словарь}\label{sec:glossary}

Текст использует немного специальных слов. Все они собраны здесь, вместе
с именем соответствующего объекта в формализации. Дальше мы стараемся
обходиться этими словами и школьной линейной алгеброй.

\begin{longtable}{@{}p{0.27\textwidth}p{0.50\textwidth}p{0.19\textwidth}@{}}
\textbf{Слово} & \textbf{Смысл} & \textbf{Имя в Lean}\\
\hline
\endhead
система размера $m$ ранга $k$ & векторы $g_1,\dots,g_m\in\RR^k$ и веса
  $t_c>0$, $\sum_c t_c=1$, с условием полноты $\sum_c t_c\,g_cg_c\T=\Id_k$
  & \code{WeightedDesign m k}\\[2pt]
сумма по подмножеству & $S_C=\sum_{c\in C}g_cg_c\T$ & \code{subsetSum}\\[2pt]
подходящее подмножество & $C$ с $S_C\succeq\Id$; строго подходящее: $S_C\succ\Id$
  & \code{Dominates}\\[2pt]
утверждение $\GTZ(m,k)$ & у всякой системы размера $m$ ранга $k$ есть
  подходящее $C$ из $k$ индексов & \code{GtzWeighted m k}\\[2pt]
$\GTZ(\cdot,k)$ & $\GTZ(m,k)$ при всех $m$ & \code{GtzWeightedAll k}\\[2pt]
граничная система & есть подходящая $k$-ка, но нет строго подходящей
  & \code{IsTie}\\[2pt]
коллинеарная пара & $g_b=\rho\,g_a$ при $a\ne b$ & \code{HasParallelPair}\\[2pt]
лемма о коллинеарной паре при $(m,k)$ & всякая граничная система размера
  $m$ ранга $k$ имеет коллинеарную пару & \code{HingeHoldsAtSize m k}\\[2pt]
длина вектора & $\ell_c=\|g_c\|^2$ (квадрат евклидовой длины); вектор
  \emph{длинный}, если $\ell_c\ge1$, \emph{строго длинный}, если $\ell_c>1$
  & \code{leverageOf}\\[2pt]
парный минор & $q_{ab}=(\ell_a-1)(\ell_b-1)-\langle g_a,g_b\rangle^2$; пара
  \emph{допустима}, если $q_{ab}>0$ & \code{pairGapMinor}\\[2pt]
определитель тройки & $D_{abc}=\det(K_{abc}-\Id)$, где $K_{abc}$ --- матрица
  Грама трёх векторов & \code{tripleGapDet}\\[2pt]
площадь пары & $w_{ab}=\ell_a\ell_b-\langle g_a,g_b\rangle^2=\|g_a\times g_b\|^2$
  & \code{crossNormSq}\\[2pt]
объём тройки & $[abc]=\det(g_a,g_b,g_c)$ & \code{tripleBracket}\\[2pt]
зазор тройки & $G_C=S_C-\Id$; у подходящей тройки $G_C\succeq0$, у граничной
  системы он вырожден & ---\\[2pt]
нулевой вектор зазора & $w\ne0$ с $S_Cw=w$; его \emph{показания} ---
  числа $\langle g_c,w\rangle$, $c\in C$ & ---\\[2pt]
напряжение & ненулевой набор $\lambda_c$ с $\sum_c\lambda_c\,g_cg_c\T=0$;
  система \emph{свободна}, если напряжений нет & ---\\[2pt]
случай ранга один («угол») & зазор $G_C=\lambda uu\T$, $\|u\|=1$, $\lambda>0$
  & ---\\[2pt]
случай ранга два & ядро зазора $G_C$ --- прямая & ---\\
\end{longtable}

Слово «подходящее» заменяет термин «доминирующее», «граничная» ---
«tie», «лемма о коллинеарной паре» --- «hinge». Буква $\Sk_k$ обозначает
пространство симметрических матриц порядка $k$; $\dim\Sk_k=k(k+1)/2$.

% ==========================================================================
\section{Постановка}

Пусть $(g,t)$ --- система размера $m$ ранга $k$. Условие полноты
\begin{equation}\label{eq:parseval}
  \sum_{c=1}^{m} t_c\, g_cg_c\T=\Id_k
\end{equation}
влечёт, что векторы $g_c$ порождают $\RR^k$, и что $t_c<1$ при $m\ge2$.
Веса в $S_C$ не входят: они связывают систему условием~\eqref{eq:parseval},
а сумма $S_C$ их не видит.

\paragraph{Исходная форма.}
Положим $v_c=\sqrt{t_c}\,g_c$ и соберём строки $v_c\T$ в матрицу $V$
размера $m\times k$. Условие~\eqref{eq:parseval} есть $V\T V=\Id_k$:
столбцы $V$ ортонормированы. Для подматрицы $B$ из строк $C$ имеем
$B\T B=\sum_{c\in C}t_c\,g_cg_c\T$. При равных весах $t_c=1/m$ условие
$S_C\succeq\Id$ превращается в $B\T B\succeq\frac1m\Id$, то есть
$\|B^{-1}\|_2\le\sqrt m$. Так что $\GTZ(m,k)$ при равных весах --- это
утверждение: у всякой $m\times k$ матрицы с ортонормированными столбцами
есть $k$ строк, образующих подматрицу $B$ с $\|B^{-1}\|_2\le\sqrt m$
(\code{GtzOriginal}). Задача восходит к работе~\cite{gtz}. Взвешенная форма
сильнее, и именно она оказывается точной (пример~\ref{ex:diamond}).

Всюду далее $s_c=t_c\ell_c=\|v_c\|^2$; след равенства~\eqref{eq:parseval}
даёт $\sum_c s_c=k$.

\begin{proposition}\label{prop:rankone}
$\GTZ(m,1)$ верно при всех $m$.
\end{proposition}
\begin{proof}
Здесь $g_c$ --- числа, $\sum_c t_cg_c^2=1=\sum_c t_c$, поэтому $g_a^2\ge1$
хотя бы при одном $a$, и годится $C=\{a\}$.
\end{proof}

% ==========================================================================
\section{Что доказано для всех рангов}

\begin{lemma}\label{lem:identity}
Положим $W=\sum_{c}(1-t_c)\,g_cg_c\T$. Для любого $C$
\[
  S_C-\Id_k=W-\sum_{c\notin C}g_cg_c\T .
\]
\end{lemma}
\begin{proof}
Вычтем~\eqref{eq:parseval} из определения $S_C$ и разобьём сумму по $C$
и дополнению: $S_C-\Id=\sum_{c\in C}(1-t_c)g_cg_c\T-\sum_{c\notin C}t_cg_cg_c\T$.
Первая сумма равна $W-\sum_{c\notin C}(1-t_c)g_cg_c\T$, и слагаемые с $t_c$
уничтожаются.
\end{proof}

\begin{corollary}\label{cor:dual}
$S_C\succeq\Id$ тогда и только тогда, когда $\sum_{c\notin C}g_cg_c\T\preceq W$.
\end{corollary}

Условие зависит только от дополнения к $C$ и от одной матрицы $W$, общей
для всех $C$. При $m\ge2$ матрица $W$ положительно определена. Возьмём $R$
с $R\T WR=\Id$ и положим $u_c=R\T g_c$; тогда
\begin{equation}\label{eq:whitened}
  S_C\succeq\Id \iff \sum_{c\notin C}u_cu_c\T\preceq\Id,
  \qquad\text{причём}\qquad \sum_c(1-t_c)\,u_cu_c\T=\Id .
\end{equation}

\begin{theorem}[\code{InductionStep.corankFloor\_holds}]\label{thm:corank}
$\GTZ(m,k)$ верно при $m\le k+2$ и любом $k$.
\end{theorem}
\begin{proof}
При $m=k$ берём все индексы: дополнение пусто.
При $m=k+1$ дополнение --- один индекс $a$, и нужно $\|u_a\|^2\le1$.
След правого равенства в~\eqref{eq:whitened} даёт $\sum_c(1-t_c)\|u_c\|^2=k$,
а $\sum_c(1-t_c)=m-1=k$. Значит $\sum_c(1-t_c)(\|u_c\|^2-1)=0$, и хотя бы
одно слагаемое неположительно.
При $m=k+2$ дополнение --- пара $\{a,b\}$, и нужно
$u_au_a\T+u_bu_b\T\preceq\Id$, то есть $\lambda_{\max}$ матрицы Грама пары
не больше единицы. Задача стала двумерной; построение Наймарка переводит
её в отбор пары в системе ранга~2, и работает теорема~\ref{thm:ranktwo}.
\end{proof}

\begin{theorem}[\code{gtz\_rank\_two}, \cite{airi}]\label{thm:ranktwo}
$\GTZ(m,k)$ верно при $k\le2$ и всех $m$.
\end{theorem}

В формализации случай $k=2$ доказан индукцией по $m$: если какой-то вектор
не длинный, он удаляется и работает предположение индукции; если все
длинные, подходящая пара строится явно, без спектральной теоремы.

\begin{theorem}[\code{not\_complexGtzWeighted\_of\_rank\_add\_two\_le\_size}]
\label{thm:complex}
Над $\CC$ (эрмитово сопряжение вместо транспонирования) утверждение
$\GTZ_\CC(m,k)$ при $k\ge2$ верно тогда и только тогда, когда $m\le k+1$.
\end{theorem}

Контрпример при $(4,2)$: четыре равноугольные прямые в $\CC^2$
(\S\ref{sec:field}). Над $\CC$ выживает ровно та область, которую над
$\RR$ теорема~\ref{thm:corank} закрывает без своего последнего шага.
Поэтому всякое доказательство вещественного случая при $m\ge k+2$
обязано где-то использовать вещественность.

\begin{example}[\code{diamondDesign}; ромб]\label{ex:diamond}
Граф $K_4$ без одного ребра, на рёбрах проводимости $c=(2,3,3,3,3)$;
матрица Кирхгофа
\[
  L=\sum_{e}c_e\,b_eb_e\T=\begin{pmatrix}8&-2&-3\\-2&8&-3\\-3&-3&6\end{pmatrix},
  \qquad\det L=180,
\]
где $b_e$ --- столбцы матрицы инцидентности. Положим $g_e=\sqrt{5c_e}\,R\T b_e$
с $R\T LR=\Id_3$ и $t_e=1/5$. Тогда $\ell_0=2$, $\ell_1=\dots=\ell_4=13/4$,
$\sum_e t_e\ell_e=3$. Для $C=\{0,1,2\}$ квадратичная форма зазора равна
$\tfrac12[(8y_1-2y_2-3y_3)^2+9y_3^2]$: она неотрицательна и обращается в
нуль при $y=(1,4,0)$. Значит $S_C\succeq\Id$, но не $S_C\succ\Id$.
Из десяти троек восемь подходят нестрого и ни одна строго; коллинеарных пар
нет. Это граничная система размера~5 ранга~3 без коллинеарных пар
(\code{not\_hingeHoldsAtSize\_five\_three}).
\end{example}

\begin{remark}[постоянная $\sqrt m$ точна]
У ромба веса равны, так что это матрица $V$ размера $5\times3$ с
ортонормированными столбцами, у которой \emph{каждая} подматрица $B$ из трёх
строк имеет $\|B^{-1}\|_2\ge\sqrt5$, и для восьми из десяти --- ровно
$\sqrt5$. Усилить единицу в $S_C\succeq\Id$, то есть $\sqrt m$ в исходной
форме, нельзя уже при $m=5$, $k=3$.
\end{remark}

% ==========================================================================
\section{Одна клетка}

\subsection{Почему $(6,3)$}

Теорема~\ref{thm:ranktwo} требует $k\ge3$, теорема~\ref{thm:corank} ---
$m\ge k+3$. Наименьшая пара, удовлетворяющая обоим условиям, есть $m=6$,
$k=3$. При $k=3$ это ещё и \emph{единственная} неизвестная клетка: всё
остальное сводится к ней (\S\ref{sub:cell}).

Число $6$ выделено: при $m=6$, $k=3$ число матриц $g_cg_c\T$ равно
$\dim\Sk_3=6$, и система~\eqref{eq:parseval} квадратна.

\begin{theorem}[\code{sizeGap\_at\_rank\_three}]
Пять векторов в $\RR^3$ всегда лежат на общей конике: есть ненулевая
$F\in\Sk_3$ с $g_c\T Fg_c=0$ при всех $c$. Для семи векторов всегда есть
напряжение. Шесть --- единственный размер, при котором не обязательны
ни коника, ни напряжение.
\end{theorem}

\subsection{Клетка $(6,3)$ решает весь ранг~3}\label{sub:cell}

\begin{proposition}[монотонность по размеру; \code{gtzWeighted\_of\_le}]
\label{prop:mono}
Из $\GTZ(m,k)$ следует $\GTZ(m',k)$ при $m'\le m$.
\end{proposition}
\begin{proof}
Достаточно $m'=m-1$. Систему размера $m-1$ дополним, расщепив один вектор
$g_1$ с весом $t_1$ на две копии с весами $t_1/2$. Условие~\eqref{eq:parseval}
сохранится. Подходящее $k$-подмножество новой системы не может содержать
обе копии: тогда $S_C$ порождалась бы $k-1$ направлениями и имела бы ранг
$<k$, а $S_C\succeq\Id$ требует ранга $k$. Значит оно даёт подходящее
$k$-подмножество исходной системы с той же $S_C$.
\end{proof}

\begin{proposition}[прогулка по напряжению; \code{gtzWeighted\_of\_forall\_smaller}]
\label{prop:walk}
Пусть $m>k(k+1)/2$. Если $\GTZ(m',k)$ верно при всех $m'<m$, то верно и
$\GTZ(m,k)$.
\end{proposition}
\begin{proof}
Матриц $g_cg_c\T$ больше, чем $\dim\Sk_k$, так что есть напряжение
$\sum_c\lambda_c g_cg_c\T=0$. Сдвигаем веса $t_c\mapsto t_c+s\lambda_c$:
условие~\eqref{eq:parseval} сохраняется при всех $s$, а сумма весов меняется
линейно. Выберем знак $s$ так, чтобы сумма весов не росла, и увеличиваем
$|s|$ до первого момента, когда какой-то вес обращается в нуль. Получили
систему с $m-1$ векторами, весами $t'_c>0$ и суммой весов $\sigma\le1$.
Положим $t''_c=t'_c/\sigma$, $g''_c=\sqrt{\sigma}\,g_c$: это система размера
$m-1$. Её подходящее $C$ даёт $\sigma S_C\succeq\Id$, откуда
$S_C\succeq\sigma^{-1}\Id\succeq\Id$. Сжатие векторов только мешает отбору,
поэтому $C$ подходит и для исходной системы.
\end{proof}

\begin{theorem}[\code{gtzWeightedAll\_of\_veroneseTop}]\label{thm:top}
При любом $k$ из $\GTZ\bigl(k(k+1)/2,\,k\bigr)$ следует $\GTZ(m,k)$ при всех $m$.
В частности $\GTZ(6,3)$ влечёт $\GTZ(m,3)$ при всех $m$.
\end{theorem}
\begin{proof}
Размеры $m\le k(k+1)/2$ --- предложение~\ref{prop:mono}, размеры больше ---
предложение~\ref{prop:walk} и индукция.
\end{proof}

Итак, на каждый ранг приходится ровно одна решающая клетка:
$m=\dim\Sk_k=k(k+1)/2$. При $k=3$ это $(6,3)$, при $k=4$ --- $(10,4)$.

\subsection{Лемма о коллинеарной паре даёт клетку}

\begin{theorem}[\code{gtzWeighted\_six\_three\_of\_hinge}]\label{thm:hinge}
Если всякая граничная система размера~6 ранга~3 имеет коллинеарную пару,
то $\GTZ(6,3)$ верно, а с ним (теорема~\ref{thm:top}) и $\GTZ(m,3)$ при всех $m$
(\code{gtzWeightedAll\_three\_of\_hinge}).
\end{theorem}
\begin{proof}[Набросок]
Пусть $D$ --- система размера~6 ранга~3.
\emph{Если в $D$ есть коллинеарная пара} $g_b=\rho g_a$, склеим её в один
вектор: $g'=\mu g_a$ с весом $t_a+t_b$, где $(t_a+t_b)\mu^2=t_a+t_b\rho^2$.
Получилась система размера~5, и по теореме~\ref{thm:corank} у неё есть
подходящая тройка. Если она не содержит $g'$, она подходит и для $D$.
Если содержит, то $\mu^2$ есть среднее чисел $1$ и $\rho^2$, поэтому
$\max(1,\rho^2)\ge\mu^2$, и тройка с $g_a$ или с $g_b$ на месте $g'$ подходит
для $D$ (\code{exists\_dominating\_triple\_of\_hasParallelPair\_sixThree}).
\emph{Если коллинеарных пар нет}, соединим $D$ непрерывным путём с фиксированной
системой-якорем, у которой есть строго подходящая тройка, --- с шестью
диагоналями икосаэдра (\code{icosaDesign\_strictly\_dominates}); путь можно
провести так, чтобы все системы на нём были без коллинеарных пар
(\code{parallelFreeReachesAnchor\_six\_three}). Множество систем со строго
подходящей тройкой открыто. Если $D$ ему не принадлежит, на пути есть
первая точка, где строгой тройки уже нет; по непрерывности у неё есть
нестрого подходящая тройка, то есть это граничная система без коллинеарных
пар --- вопреки предположению. Значит у $D$ есть строго подходящая тройка.
\end{proof}

Это единственный недостающий шаг. При $m=5$ соответствующее утверждение
ложно (ромб); разница между пятью и шестью векторами не количественная:
пять обязаны лежать на конике, шесть --- нет.

\subsection{Эквивалентные формы остатка}

Формализация содержит реестр из семи аксиом. Пять из них --- утверждения о
классах систем размера~6 ранга~3, и доказано, что их конъюнкция
\emph{равносильна} лемме о коллинеарной паре при $(6,3)$
(\code{onPathRegistryCollapse}, \code{allFiveOnPath\_iff\_hingeHoldsAtSize\_six\_three}).
Две оставшиеся относятся к рангам $\ge4$ (\S\ref{sec:allranks}).

Два сведения полезны на практике.
\begin{itemize}
\item \emph{Только свободные системы.} Всякая граничная система размера~6
с напряжением имеет коллинеарную пару (\code{sixThree\_stress\_trichotomy}
и доказанные ветви). Свободная система коллинеарной пары иметь не может
(коллинеарная пара --- это напряжение из двух слагаемых). Поэтому лемма
равносильна утверждению: \emph{граничной системы размера~6 ранга~3 со
свободными матрицами $g_cg_c\T$ не существует}
(\code{StressFreeHingeHoldsSixThree}). В свободной системе веса однозначно
восстанавливаются по направлениям: система~\eqref{eq:parseval} квадратна и
невырождена.
\item \emph{Три неравенства.} Лемма равносильна системе из трёх
полиномиальных неравенств на элементы проектора $VV\T$ и веса, требуемой
при каждой системе без коллинеарных пар
(\code{offsetDominatesSomewhere\_iff\_hingeHoldsAtSize\_six\_three}).
\end{itemize}

Для ранга~3 этого достаточно. Место клетки $(6,3)$ во всей задаче описано
в~\S\ref{sec:allranks}: при двух общих предположениях о высших рангах
утверждение для всех рангов \emph{равносильно} $\GTZ(6,3)$
(\code{everyRank\_iff\_sixThree}).

% ==========================================================================
\section{Анатомия граничной системы размера шесть}\label{sec:anatomy}

Всюду в этом разделе $D$ --- граничная система размера~6 ранга~3, и мы
хотим найти в ней коллинеарную пару. Пусть $C=\{a,b,c\}$ --- подходящая
тройка: $S_C\succeq\Id$, но ни одна тройка не подходит строго.

\subsection{Критерий через матрицу Грама}

Пусть $G=[g_a\ g_b\ g_c]$ --- матрица $3\times3$ из столбцов тройки. Тогда
$S_C=GG\T$, а $K_C=G\T G$ --- матрица Грама тройки. Матрицы $GG\T$ и $G\T G$
имеют одинаковые собственные значения, поэтому
\[
  S_C\succeq\Id\iff K_C\succeq\Id,\qquad S_C\succ\Id\iff K_C\succ\Id .
\]
Так условие на тройку становится условием на шесть чисел:
$\ell_a,\ell_b,\ell_c$ и три попарных произведения. По критерию Сильвестра
$K_C-\Id\succ0$ тогда и только тогда, когда
\begin{equation}\label{eq:sylvester}
  \ell_a>1,\qquad q_{ab}>0,\qquad D_{abc}>0,
\end{equation}
где $q_{ab}$ --- парный минор, $D_{abc}=\det(K_C-\Id)$
(\code{tripleGram\_posDef\_iff\_pairVocabulary}). Первое условие видит один
вектор, второе --- пару, и только третье --- тройку. Значит в граничной
системе для \emph{каждой} тройки нарушено одно из трёх: какой-то её вектор
не строго длинный, какая-то пара недопустима, или определитель тройки
неположителен (\code{isTie\_triple\_trichotomy}). Подчеркнём: условие
$D_{abc}\le0$ появляется только у троек, все пары которых допустимы.

\subsection{Все векторы длинные}

\begin{proposition}[\code{leverage\_one\_le\_of\_isTie}]\label{prop:long}
В граничной системе размера~6 ранга~3 все шесть векторов длинные: $\ell_c\ge1$.
\end{proposition}
\begin{proof}
Пусть $\ell_a<1$. Матрица $M=\Id-t_ag_ag_a\T=\sum_{c\ne a}t_cg_cg_c\T$
положительно определена. Положим $h_c=\sqrt{1-t_a}\,M^{-1/2}g_c$ и
$t'_c=t_c/(1-t_a)$ для пяти $c\ne a$: это система размера~5, и по
теореме~\ref{thm:corank} у неё есть подходящая тройка $T$, то есть
$(1-t_a)M^{-1/2}S_TM^{-1/2}\succeq\Id$, или
\begin{equation}\label{eq:funnel}
  S_T\;\succeq\;\frac{\Id-t_ag_ag_a\T}{1-t_a}.
\end{equation}
Собственные значения правой части: $1/(1-t_a)>1$ на $g_a^\perp$ и
$(1-t_a\ell_a)/(1-t_a)$ вдоль $g_a$; последнее больше единицы ровно при
$\ell_a<1$. Тогда $S_T\succ\Id$ --- вопреки граничности.
\end{proof}

\subsection{Воронка: вектор длины ровно один}\label{sub:funnel}

\begin{proposition}[\code{isTie\_sixThree\_allHeavy\_or\_funnel},
\code{unitAtom\_parallel\_of\_two\_readings\_zero}]\label{prop:funnel}
Пусть в граничной системе $D$ размера~6 ранга~3 вектор $g_a$ имеет $\ell_a=1$.
Тогда есть подходящая тройка $T\not\ni a$, для которой
\[
  S_Tg_a=g_a,\qquad \sum_{c\in T}\langle g_c,g_a\rangle^2=1 .
\]
Если при этом два из трёх показаний $\langle g_c,g_a\rangle$, $c\in T$, равны
нулю, то в $D$ есть коллинеарная пара: вектор $g_a$ кратен третьему вектору
тройки.
\end{proposition}
\begin{proof}
При $\ell_a=1$ правая часть~\eqref{eq:funnel} равна $\Id+\mu P$, где
$\mu=t_a/(1-t_a)>0$ и $P=\Id-g_ag_a\T$ --- проектор на $g_a^\perp$. Значит
зазор $G_T=S_T-\Id\succeq\mu P$ положительно определён на $g_a^\perp$. Он
вырожден (система граничная), и любой его нулевой вектор $w$ даёт
$0=w\T G_Tw\ge\mu\|Pw\|^2$, то есть $w\parallel g_a$. Отсюда $S_Tg_a=g_a$,
то есть $g_a=\sum_{c\in T}\langle g_c,g_a\rangle g_c$; скалярно умножая на
$g_a$, получаем сумму квадратов показаний, равную $\|g_a\|^2=1$. Если два
показания нулевые, $g_a$ кратен третьему вектору.
\end{proof}

Так что граничная система размера~6 либо состоит из строго длинных
векторов, либо «воронка»: вектор длины один воспроизводится подходящей
тройкой, его не содержащей. Ветвь воронки разобрана почти до конца, и
разбор сводится к одному утверждению о размере~5.

\begin{proposition}[ветвь воронки]\label{prop:funnelbranch}
Пусть $D$ --- граничная система размера~6 ранга~3 без коллинеарных пар,
$\ell_a=1$, и пусть $\mathcal T$ --- множество троек $T\not\ni a$,
удовлетворяющих~\eqref{eq:funnel}; это в точности подходящие тройки
укороченной системы размера~5 из доказательства
предложения~\ref{prop:long}. Тогда:
\begin{enumerate}
\item укороченная система граничная, и $\mathcal T\ne\varnothing$;
\item каждая $T\in\mathcal T$ воспроизводит $g_a$: $S_Tg_a=g_a$
  (\code{dominates\_and\_fixes\_of\_deflatedGapBound});
\item ни одна $T\in\mathcal T$ не является углом: её $e_2=q_{xy}+q_{xz}+q_{yz}
  \ge\bigl(t_a/(1-t_a)\bigr)^2>0$, и оценка точна
  (\code{pairMinorTotal\_floor\_of\_deflatedGapBound});
\item если $T,T'\in\mathcal T$ различаются одним обменом, то заменяемый
  вектор слеп к $g_a$, и у $T$ не может быть двух таких соседей; формально это
  \code{unitAtom\_two\_mirror\_swaps\_absurd} для обменов у разных членов $T$ и
  \code{unitAtom\_second\_swap\_absurd} для двух обменов у одного члена;
\item тройки пятиэлементного множества --- это дополнения рёбер графа $K_5$,
  обмен --- смежность рёбер, и семейство рёбер, в котором каждое ребро смежно
  не более чем с одним другим, содержит не более трёх рёбер
  (\code{swapDegree\_le\_one\_card\_le\_three}); значит $|\mathcal T|\le3$;
\item вне каждой $T\in\mathcal T$ есть строго длинный вектор $u$ с
  $q_{au}<-1$ (\code{isTie\_sixThree\_funnel\_structure}).
\end{enumerate}
\end{proposition}

Пункт~2: зазор $G_T$ вырожден и положительно определён на $g_a^\perp$, как
в доказательстве предложения~\ref{prop:funnel}. Пункт~3: $e_2$ --- квадратичная
форма на матрицах, и разложение $G_T=\mu P+R$ с $R\succeq0$ даёт
$e_2(G_T)\ge\mu^2e_2(P)=\mu^2$. Пункт~4: зеркало (\S\ref{sub:mirror})
даёт слепоту; два слепых вектора у разных членов $T$ дают коллинеарную пару
по предложению~\ref{prop:funnel}, а два обмена у одного члена делают слепыми
три вектора, и тогда условие~\eqref{eq:parseval}, прочитанное на $g_a$,
требует $1-t_a\le t_x+t_y<1-t_a$.

\emph{Следствие.} Если у всякой граничной системы размера~5 ранга~3 не
меньше четырёх подходящих троек, ветвь воронки пуста: у $D$ есть
коллинеарная пара. Численно это так с запасом: у $71$ найденной граничной
системы размера~5 подходящих троек от $6$ до $8$ (у ромба восемь).
Само утверждение «не меньше четырёх» не доказано; это
задача~\ref{task:five} ниже.

\subsection{Пять балансов}\label{sub:budgets}

Условие~\eqref{eq:parseval} даёт пять точных тождеств, верных для
\emph{всякой} системы размера $m$ ранга~3, граничной или нет:
\begin{align}
  &\sum_c t_c\ell_c=3, \label{eq:b1}\\
  &\sum_{a,b}t_at_b\langle g_a,g_b\rangle^2=3, \label{eq:b2}\\
  &\sum_{a,b}t_at_b\,w_{ab}=6, \label{eq:b3}\\
  &\sum_{a,b}t_at_b\,q_{ab}=1, \label{eq:b4}\\
  &\sum_{a<b<c}t_at_bt_c\,[abc]^2=1. \label{eq:b5}
\end{align}
Первое --- след~\eqref{eq:parseval}, второе --- квадрат нормы Фробениуса
обеих частей, третье --- разность $(\sum t_c\ell_c)^2$ и второго, пятое ---
формула Коши--Бине для $\det(V\T V)=1$. Четвёртое
(\code{pairMinor\_budget}) получается подстановкой
$q_{ab}=w_{ab}-\ell_a-\ell_b+1$: $6-3-3+1=1$. Его следствие: диагональ
$\sum_a t_a^2q_{aa}=\sum_at_a^2(1-2\ell_a)$ отрицательна, так что
внедиагональная часть положительна, и \emph{в любой системе есть допустимая
пара} (\code{exists\_offDiag\_pairMinor\_pos}); коллинеарная пара, напротив,
вносит в~\eqref{eq:b4} сильно отрицательный вклад $q_{ab}=-(\ell_a+\ell_b-1)$.

Ещё одно тождество: взвешенная сумма определителей всех троек
$\sum_{a,b,c}t_at_bt_cD_{abc}=-4$ одна и та же у всех систем
(\code{weighted\_aggregate\_design\_blind}). Поэтому никакой аргумент,
использующий только усреднённые по тройкам величины, закрыть клетку не
может: он не отличает одну систему от другой.

\subsection{Закон присоединённой матрицы}\label{sub:adj}

Это единственное новое тождество раздела, и оно связывает две ветви
разбора.

\begin{proposition}[\code{pairGapMinor\_eq\_pairMinorTotal\_mul\_reading\_first}
и парные к ней]\label{prop:adj}
Пусть тройка $\{a,b,c\}$ подходит, $w$ --- единичный нулевой вектор её
зазора ($S_Cw=w$, $\|w\|=1$), $s=(\langle g_a,w\rangle,\langle g_b,w\rangle,
\langle g_c,w\rangle)$ --- показания, $e_2=q_{ab}+q_{ac}+q_{bc}$. Тогда
\[
  q_{bc}=e_2\,s_a^2,\qquad q_{ac}=e_2\,s_b^2,\qquad q_{ab}=e_2\,s_c^2,
  \qquad s_a^2+s_b^2+s_c^2=1 .
\]
\end{proposition}
\begin{proof}
Пусть $N=K_C-\Id$. Координата $a$ вектора $Ns$ равна
$\langle g_a,\sum_{c}\langle g_c,w\rangle g_c\rangle-\langle g_a,w\rangle
=\langle g_a,S_Cw-w\rangle=0$; значит $Ns=0$. Далее $\|s\|^2=w\T S_Cw=\|w\|^2=1$.
Присоединённая матрица $\adj N$ симметрична, $N\cdot\adj N=\det N\cdot\Id=0$,
и её ранг не больше единицы, так как $N$ вырождена. Если $\rank N=2$,
столбцы $\adj N$ лежат в $\ker N=\RR s$, откуда $\adj N=\mu ss\T$, а след даёт
$\mu=\tr\adj N=e_2$. Если $\rank N\le1$, то $\adj N=0$ и $e_2=0$. В обоих
случаях $\adj N=e_2\,ss\T$, и диагональ этого равенства --- утверждение.
\end{proof}

\begin{corollary}\label{cor:adj}
\begin{enumerate}
\item В случае ранга один ($e_2=0$) все три парных минора подходящей тройки
равны нулю: тройка «удаляет свой треугольник» из графа допустимых пар.
\item В случае ранга два ($e_2>0$) парный минор $q_{bc}$ равен нулю тогда
и только тогда, когда третий вектор $g_a$ «слеп» к нулевому вектору:
$\langle g_a,w\rangle=0$.
\item Все пары внутри подходящей тройки допустимы нестрого, $q\ge0$
(\code{pairGapMinor\_nonneg\_of\_dominates}); значит \emph{недопустимая пара
($q<0$) всегда имеет участника вне каждой подходящей тройки}
(\code{inadmissible\_pair\_not\_inside\_dominates}).
\end{enumerate}
\end{corollary}

Нормированные парные миноры подходящей тройки --- это распределение
вероятностей, и оно равно квадратам показаний нулевого вектора. У ромба
(пример~\ref{ex:diamond}) для тройки $\{0,1,2\}$ показания в квадрате
равны $(0.75,0.125,0.125)$, а миноры $(4.5,0.75,0.75)$ с $e_2=6$.

То же рассуждение в исходном пространстве даёт закон четвёрок. Зазор
$G=S_C-\Id$ симметричен, $Gw=0$, $\|w\|=1$, и по тем же причинам
$\adj G=e_2(G)\,ww\T$, где $e_2(G)$ --- сумма главных миноров второго порядка
матрицы $G$, равная $e_2$ (\code{secondInvariantOfThree\_gap\_eq\_pairMinorTotal}).
По формуле для определителя возмущения ранга один,
\begin{equation}\label{eq:fourset}
  \det\bigl(G+vv\T\bigr)=\det G+v\T(\adj G)v=e_2\,\langle v,w\rangle^2
  \qquad\text{для всех }v\in\RR^3
\end{equation}
(\code{det\_add\_atomMatrix\_of\_unit\_null}). Поэтому четвёрка $C\cup\{d\}$
подходит строго тогда и только тогда, когда $e_2>0$ и $\langle g_d,w\rangle\ne0$
(\code{fourSet\_posDef\_of\_reading\_ne\_zero}). Условие~\eqref{eq:parseval}
заставляет какой-то вектор вне $C$ читать $w$ ненулевым показанием
(\code{OutsideReadingFloor}). Итог: \emph{в случае ранга два подходящая
тройка всегда продолжается до строго подходящей четвёрки, а угол --- никогда}
($e_2=0$ обращает~\eqref{eq:fourset} в нуль). Четвёрки разделяют две ветви
\S\ref{sub:split}.

\subsection{Зеркало: два обмена на одном нулевом векторе}\label{sub:mirror}

\begin{proposition}[\code{mirror\_reading\_smul\_eq},
\code{hasParallelPair\_of\_mirror\_nullDominators}]\label{prop:mirror}
Пусть тройки $C\ni p$ и $C'=C\setminus\{p\}\cup\{q\}$ обе подходят и имеют
общий нулевой вектор $w$. Тогда
$\langle g_p,w\rangle\,g_p=\langle g_q,w\rangle\,g_q$. В частности, если
$\langle g_p,w\rangle\ne0$, векторы $g_p,g_q$ коллинеарны.
\end{proposition}
\begin{proof}
$(S_C-S_{C'})w=g_p\langle g_p,w\rangle-g_q\langle g_q,w\rangle=0$.
\end{proof}

Вместе с предложением~\ref{prop:adj}: в системе без коллинеарных пар
показание $\langle g_p,w\rangle$ обязано быть нулевым, то есть парный минор
двух общих векторов $C\cap C'$ равен нулю
(\code{mirror\_shared\_pairGapMinor\_eq\_zero}). \emph{Каждый такой обмен
стоит системе без коллинеарных пар одного уравнения $q=0$.} Подходящая
тройка вектора длины один допускает не более одного такого обмена
(\code{unitAtom\_two\_mirror\_swaps\_absurd}).

Это объясняет ромб: у него восемь подходящих троек и восемнадцать пар
троек, различающихся одним обменом, но \emph{никакие две не имеют общего
нулевого вектора} (восемь нулевых направлений попарно различны; прямое
вычисление). Зеркало там не срабатывает ни разу, и потому ромб --- граничная
система без коллинеарных пар. Воронка (\S\ref{sub:funnel}) --- противоположный
режим: все подходящие тройки укороченной системы имеют общий нулевой
вектор $g_a$.

Нулевой вектор можно исключить из всех этих условий. Полная форма закона
присоединённой матрицы, $\adj(K_C-\Id)=e_2\,ss\T$ поэлементно, даёт для
любого $v$ равенство $e_2\,\langle v,w\rangle^2=Q_C(v)$, где $Q_C$ --- явная
квадратичная форма от трёх показаний $\langle g_a,v\rangle,\langle g_b,v\rangle,
\langle g_c,v\rangle$ (\code{ProbeFreeHingeTrigger}). Поэтому условие
зеркала --- полиномиальное уравнение на длины и попарные произведения, без
собственных векторов; для системы из шести векторов это и нужно.

\subsection{Недопустимая пара: что она даёт и чего не даёт}

Недопустимость $q_{ab}\le0$ равносильна $w_{ab}\le\ell_a+\ell_b-1$
(\code{crossNormSq\_le\_of\_pairGapMinor\_nonpos}), и эта оценка точна:
есть системы с недопустимыми длинными парами, у которых отношение
$w_{ab}/(\ell_a+\ell_b-1)$ равно $0.999999$. Коллинеарность --- это $w_{ab}=0$.
Между $q\le0$ и $w=0$ лежит зазор, который не закрывается никаким
усилением этой оценки: закрыть его может только использование граничности.

\subsection{Ранг зазора и две ветви}\label{sub:split}

Зазор $G_C=S_C-\Id$ подходящей тройки неотрицательно определён и вырожден.
Его ранг равен $0$, $1$ или $2$.

\begin{proposition}\label{prop:rank0}
В системе размера~6 ранга~3 ни одна тройка не имеет $S_C=\Id$, если система
граничная.
\end{proposition}
\begin{proof}
Пусть $S_C=\Id$, то есть $g_a,g_b,g_c$ ортонормированы. Тогда для
дополнительной тройки $\{d,e,f\}$ из~\eqref{eq:parseval}
\[
  \sum_{x\in\{d,e,f\}}t_xg_xg_x\T=\Id-\sum_{c\in C}t_cg_cg_c\T
  =\operatorname{diag}(1-t_a,1-t_b,1-t_c)
\]
в базисе $(g_a,g_b,g_c)$. Пусть $\tau=\max(t_d,t_e,t_f)$. Тогда
$S_{def}\succeq\tau^{-1}\sum_xt_xg_xg_x\T=\tau^{-1}\operatorname{diag}(1-t_c)$,
а $1-t_c>t_d+t_e+t_f\ge\tau$, так как все шесть весов положительны.
Значит $S_{def}\succ\Id$ --- строго подходящая тройка.
\end{proof}

Остаются две ветви, и они ведут себя по-разному.
\begin{itemize}
\item \textbf{Ранг один («угол»).} $S_C=\Id+\lambda uu\T$, $\|u\|=1$,
$\lambda>0$. Ожидается, что таких граничных систем над $\RR$ нет вовсе.
Это задача о \emph{пустоте}.
\item \textbf{Ранг два.} Ядро зазора --- прямая $\RR w$. Такие граничные
системы существуют (ромб при $m=5$); нужно доказать, что при $m=6$ каждая
из них имеет коллинеарную пару. Это задача о \emph{жёсткости}.
\end{itemize}

% ==========================================================================
\section{Ветвь ранга один}\label{sec:corner}

Пусть $S_C=\Id+\lambda uu\T$ для тройки $C=\{a,b,c\}$.

\subsection{Что известно точно}

Матрица Грама $K_C$ имеет собственные значения $1,1,1+\lambda$, и
$K_C-\Id=\lambda vv\T$ с единичным $v=G\T u/\sqrt{1+\lambda}$, то есть
$v_e=\langle g_e,u\rangle/\sqrt{1+\lambda}$. Отсюда сразу
(\code{CornerAxisCalculus}):
\begin{gather}
  \lambda\,\langle u,g_e\rangle^2=(1+\lambda)(\ell_e-1)\quad(e\in C),
  \qquad \sum_{e\in C}(\ell_e-1)=\lambda,\qquad
  \sum_{e\in C}\langle u,g_e\rangle^2=1+\lambda, \label{eq:cornerread}\\
  q_{ab}=q_{ac}=q_{bc}=0 . \label{eq:cornerminors}
\end{gather}
Показания оси $u$ внутри тройки не оцениваются, а \emph{определяются}
длинами; вектор тройки имеет длину ровно один тогда и только тогда, когда он
ортогонален оси. Все пары внутри тройки лежат на границе допустимости, и
остаётся опровергнуть ровно десять троек: $\binom63$ минус десять троек с
двумя и тремя векторами из $C$, которые отпадают уже
по~\eqref{eq:sylvester}--\eqref{eq:cornerminors}
(\code{corner\_posDef\_triple\_inter\_le\_one}).

\subsection{Разбор по числу единичных векторов тройки}

По~\eqref{eq:cornerread} число векторов $C$ с $\ell_e=1$ равно числу
ортогональных оси.
\begin{itemize}
\item \emph{Три}: невозможно, так как $\sum_e\langle u,g_e\rangle^2=1+\lambda>0$.
\item \emph{Два}: тогда третий вектор $g_c$ коллинеарен оси,
$\ell_c=1+\lambda$, а $g_a,g_b$ --- ортонормированный базис плоскости
$u^\perp$. Граничной системы с такой тройкой нет
(\code{corner\_twoAxisZero\_absurd}).
\item \emph{Один}: единичный вектор тройки --- это вектор длины один, так что
весь случай лежит в ветви воронки (\S\ref{sub:funnel}) и закрывается вместе
с ней задачей~\ref{task:five}. Независимо от этого он разобран прямо: две
«ячейки», H и B, по тому, какие четвёрки строго подходят; каждая сведена к
системе полиномиальных неравенств без обратных матриц (для ячейки B ---
семь знаков определителей троек, \code{OneAxisZeroBothLightChartSystem}),
численно несовместной, причём запас пропорционален нижней границе весов
(около $2\tau$ при $t_c\ge\tau$), так что сертификат обязан быть нестрогим.
Прямой сертификат не найден.
\item \emph{Ноль}: все три вектора строго длинные. Это единственный случай,
совместимый с ветвью, где все шесть векторов строго длинные; он сведён к
двум утверждениям E1 и E2 ниже. \emph{Открыто.}
\end{itemize}

\subsection{Два утверждения основного случая}

Пусть $d,d',d''$ --- три вектора вне $C$.

\paragraph{E1: вне тройки есть допустимая пара.}
Над $\RR$ и над $\CC$ одинаково; доказывается, когда доказывается, только
из положительной определённости и~\eqref{eq:parseval}. Инструмент ---
\emph{оценка спаривания} (\code{pairing\_cap}): для любой пары $i,j$
\begin{equation}\label{eq:cap}
  t_it_j\langle g_i,g_j\rangle^2\le(1-t_i\ell_i)(1-t_j\ell_j),
\end{equation}
потому что матрица $\Id-t_ig_ig_i\T-t_jg_jg_j\T=\sum_{c\ne i,j}t_cg_cg_c\T$
неотрицательно определена, а её определитель равен правой части минус
левая. Подстановка~\eqref{eq:cap} в парный минор даёт критерий в двух
числах на вектор: с $a_c=t_c(\ell_c-1)$, $A_c=1-t_c$,
\[
  \frac{a_i}{A_i}+\frac{a_j}{A_j}>1\ \Longrightarrow\ q_{ij}>0
  \qquad(\code{pairGapMinor\_pos\_of\_normalizedExcess}).
\]
По~\eqref{eq:b1} сумма $\sum_ca_c$ равна~2, а доля тройки
$\sum_{e\in C}a_e$ строго меньше единицы (она равна
$\frac{\lambda}{1+\lambda}\,u\T(\sum_{e\in C}t_eg_eg_e\T)u\le\frac\lambda{1+\lambda}$).
Значит три внешних вектора несут больше единицы, и E1 доказано во всех
режимах, кроме одного: все три внешних вектора длинные, а тройка несёт не
меньше половины (\code{CornerAdmissibleGateway}). В этом режиме E1 сведено к
одному полиномиальному неравенству в шести переменных, и установлено, что
\emph{сертификат вида «сумма произведений линейных ограничений»
не существует ни в какой степени}: найдено двупараметрическое семейство
точек, где цель обращается в нуль при стремлении одного веса к нулю
(\code{GatewayCertificateObstruction}). Нужен либо сертификат с квадратами,
либо рукописное доказательство. Само утверждение численно верно во всех
проверенных точках.

\paragraph{E2: какой-то вектор тройки «расплачивается» за внешнюю пару.}
Если $\{d,d'\}$ --- допустимая внешняя пара и $e\in C$ таков, что
$D_{edd'}>0$, то по~\eqref{eq:sylvester} тройка $\{e,d,d'\}$ подходит строго,
и граничность нарушена (\code{tripleGram\_posDef\_of\_admissible\_of\_repay}).
E2 утверждает, что такой $e$ найдётся всегда. Здесь и только здесь в ветви
ранга один нужна вещественность: \emph{существует комплексная граничная
система ранга один} --- шесть векторов в $\CC^3$, выполняющих~\eqref{eq:parseval}
с эрмитовым сопряжением, у которых $S_C=\Id+\lambda uu^*$ и ни одна из
двадцати троек не подходит строго. У неё E1 выполнено, а E2 нарушено во
всех девяти парах $(e,\text{внешняя пара})$. Поэтому E2 нельзя доказать
никаким рассуждением, верным и над $\CC$. Механизм различия: показания
$\langle g_e,u\rangle$ связаны с длинами унитарной группой, и $U(3)$
просторнее $O(3)$. Численно E2 над $\RR$ не нарушается, но запас не
отделён от нуля (при спуске доходит до $7\cdot10^{-6}$), так что
сертификат обязан быть нестрогим. Известна точная форма остатка: три числа
$D_{edd'}$, $e\in C$, --- диагональ явной симметрической матрицы $X$
порядка~3, у которой $\det X=\lambda\,q_{dd'}^2\,[u\,d\,d']^2$, а
положительных собственных значений не больше одного
(\code{CornerRepaymentMatrix}). Поэтому никакая взвешенная сумма по $e$ не
может быть положительной всегда: нужен выбор одного $e$, а не усреднение.
Вещественный ингредиент назван: это равенство между произведением трёх
попарных произведений и его действительной частью, которое над $\RR$
выполнено тождественно, а над $\CC$ нарушено у 19 из 20 троек
комплексной системы. \emph{Открыто.}

Если E1 и E2 доказаны, ветвь ранга один пуста в основном случае.

% ==========================================================================
\section{Ветвь ранга два}\label{sec:line}

Пусть ядро зазора --- прямая $\RR w$, $\|w\|=1$, $S_Cw=w$, и по
предложению~\ref{prop:adj} $e_2>0$. Ветвь делится по числу векторов
тройки, слепых к $w$ ($\langle g_e,w\rangle=0$); по следствию~\ref{cor:adj}
это число нулевых парных миноров внутри тройки.

\paragraph{Два слепых вектора.} Тогда $w=\langle g_c,w\rangle g_c$, откуда
$g_c\parallel w$, $\ell_c=1$ и $g_a,g_b\perp g_c$: тройка состоит из
единичного вектора и двух векторов в ортогональной плоскости. В системе
есть вектор длины один, так что случай лежит в ветви воронки и закрывается
вместе с ней (задача~\ref{task:five}). Прямое рассуждение при размере~6
даёт: либо коллинеарная пара, либо единичный вектор $g_c$ лежит в плоскости
двух векторов с нулевым парным минором, с собственными показаниями в роли
координат (\code{k2SixThree\_parallel\_or\_boundary\_plane}); дальнейшее
сведение этой конфигурации к двум плоскостям выписано, но не
формализовано. Для размера~5 тот же случай был приведён к явной области в
$\RR^8$ (все условия --- линейные формы в подходящих координатах), и эта
область \emph{доказана пустой} сертификатом из $28$ и $34$ произведений
линейных форм с коэффициентами $\tfrac12,1,2$, найденным линейным
программированием за доли секунды (\code{k2Chart\_kill}); мост от системы к
области несёт перенесённые гипотезы (\code{k2FiveTwoZero\_kill}). Этот путь
теперь запасной.

\paragraph{Один слепой вектор.} Скажем, $\langle g_a,w\rangle=0$; тогда
$w$ лежит в плоскости $g_b,g_c$ и $q_{bc}=0$. Численно случай пуст при
каждом фиксированном нижнем пороге на веса, но точка минимума --- точка
Чебышёва: шесть ограничений активны одновременно, и спуск не сходится.
Нужно решать систему активных ограничений точно, а не спускаться.
\emph{Открыто.}

\paragraph{Ни одного слепого вектора.} Все три парных минора внутри тройки
строго положительны. Это общий случай и сердцевина задачи. Разобран по
тому, есть ли в системе недопустимая пара:
\begin{itemize}
\item[(а)] \emph{Все пятнадцать пар допустимы и все векторы строго длинные.}
Тогда по~\eqref{eq:sylvester} у каждой тройки $D\le0$. Нужно показать, что
такой граничной системы нет. Подсчёт троек (теорема Мантеля для графа
допустимых пар формализована) сам по себе знака определителя не даёт:
$D\le0$ у всех троек --- это \emph{условие}, а не вывод. Нужен
аналитический аргумент. \emph{Открыто.}
\item[(б)] \emph{Есть недопустимая пара $\{a,b\}$.} По следствию~\ref{cor:adj}
она не лежит ни в одной подходящей тройке целиком. Нужно вывести
коллинеарность $g_a,g_b$. Точная оценка $w_{ab}\le\ell_a+\ell_b-1$ одна
этого не даёт (\S\ref{sec:anatomy}); недостающий ингредиент обязан
использовать граничность --- например, через зеркало
(\S\ref{sub:mirror}): нужен второй подходящий набор с общим нулевым
вектором. \emph{Открыто; это последний логический стык леммы.}
\end{itemize}

% ==========================================================================
\section{Точный список того, что осталось доказать}\label{sec:list}

Граничная система размера~6 ранга~3 либо содержит вектор длины один
(ветвь воронки), либо все её векторы строго длинные
(предложение~\ref{prop:long}). Лемма о коллинеарной паре при $(6,3)$ следует
из совокупности пяти утверждений ниже; сборка кусков в лемму --- техническая
работа, частично выполненная.

\begin{task}[ветвь воронки; утверждение о размере пять]\label{task:five}
У всякой граничной системы размера~5 ранга~3 не меньше четырёх подходящих
троек. (Численно: от шести до восьми, у ромба восемь.) По
предложению~\ref{prop:funnelbranch} это закрывает всю ветвь воронки, и с ней
ячейки H, B ветви ранга один и случай двух слепых векторов ветви ранга два.
\end{task}

Дальше все векторы строго длинные.

\begin{task}[ранг один, E1]\label{task:e1}
В режиме «три внешних вектора длинные, тройка несёт не меньше половины
избытка» вне тройки есть допустимая пара. Сведено к одному полиномиальному
неравенству; сертификат из произведений линейных ограничений невозможен
ни в какой степени, требуется сертификат с квадратами или ручное
доказательство.
\end{task}

\begin{task}[ранг один, E2]\label{task:e2}
При допустимой внешней паре $\{d,d'\}$ найдётся $e\in C$ с $D_{edd'}>0$.
Доказательство обязано использовать вещественность; сертификат обязан быть
нестрогим.
\end{task}

\begin{task}[ранг два, один слепой вектор]\label{task:k1}
Пустота; через точное решение системы активных ограничений, а не спуском.
\end{task}

\begin{task}[ранг два, ни одного слепого]\label{task:k0}
(а) Нет граничной системы, у которой все векторы строго длинные и все
пятнадцать пар допустимы. (б) Недопустимая пара граничной системы
коллинеарна.
\end{task}

Задачи~\ref{task:e1}, \ref{task:e2} вместе дают пустоту ветви ранга один
при строго длинных векторах; задачи~\ref{task:k1}, \ref{task:k0} --- жёсткость
ветви ранга два. Запасные пути (ячейки H, B напрямую; область в $\RR^8$)
остаются и могут оказаться короче задачи~\ref{task:five}.

Всё остальное в разборе --- теоремы: критерий через матрицу Грама,
длина всех векторов, воронка и её предложение~\ref{prop:funnelbranch}, пять
балансов, закон присоединённой матрицы и закон четвёрок, зеркало и его
бесспектральная форма, исключение ранга нуль, точный расчёт угла,
исключение троек с двумя векторами из $C$ и случая двух единичных векторов,
область ранга два при размере~5.

% ==========================================================================
\section{Дорога до ранга три}\label{sec:rankthree}

Она состоит из уже доказанных стрелок и одного недостающего звена:
\[
  \underbrace{\text{задачи \ref{sec:list}}}_{\text{открыто}}
  \ \Longrightarrow\ \text{лемма при }(6,3)
  \ \overset{\text{т.~\ref{thm:hinge}}}{\Longrightarrow}\ \GTZ(6,3)
  \ \overset{\text{т.~\ref{thm:top}}}{\Longrightarrow}\ \GTZ(m,3)\ \forall m
  \ \Longrightarrow\ \text{исходная форма при }k=3 .
\]
Иначе говоря, для ранга~3 нужно доказать ровно лемму о коллинеарной паре
при $(6,3)$, то есть закрыть список~\S\ref{sec:list}. Заметим: годится и
любое прямое доказательство $\GTZ(6,3)$, минуя лемму (например, построение
строго подходящей тройки у каждой неграничной системы и разбор граничных);
теорема~\ref{thm:top} доводит его до всех размеров.

Три рабочих соображения.
\begin{enumerate}
\item Достаточно рассматривать свободные системы: веса тогда --- рациональные
функции направлений, и все балансы~\S\ref{sub:budgets} становятся тождествами
в одних направлениях.
\item Вещественность нужна только в E2 (и, возможно, в
задаче~\ref{task:k0}); всё остальное в списке --- «бесполевые» утверждения,
которые можно искать методами положительной определённости.
Задача~\ref{task:five} и вовсе про размер~5, где всё известно.
\item Любой сертификат обязан быть нестрогим: множество граничных систем ---
это множество \emph{минимумов} функции
$\Phi=\max_{|C|=3}\lambda_{\min}(S_C)$, а не её линия уровня, и запасы на
нём не отделены от нуля.
\end{enumerate}

% ==========================================================================
\section{Дорога до произвольного ранга}\label{sec:allranks}

\subsection{Скелет индукции}

Формализация содержит полный вывод $\GTZ(m,k)$ при \emph{всех} $m$ и $k$ из
семи аксиом реестра (\code{skeletonGtzWeightedAllEveryRank}); пять из них ---
клетка $(6,3)$ (\S\ref{sub:cell}), две --- утверждения о высших рангах. Вот
его устройство. Индукция по рангу $k$; при $k\le2$ всё доказано. Пусть $k\ge3$
и ранг $k-1$ закрыт полностью. Размеры при ранге $k$ делятся на четыре зоны.
\begin{enumerate}
\item $m\le k+2$: теорема~\ref{thm:corank}.
\item $k+3\le m\le2k-1$: сводится к рангу $k-1$. К системе размера $m$ ранга
$k$ есть \emph{дополнительная} система того же размера и ранга $m-k$ с теми
же весами, в которой подходящие $k$-подмножества первой --- в точности
дополнения подходящих $(m-k)$-подмножеств второй
(\code{weighted\_naimark\_duality}). При $m=2k-1$ ранг дополнения равен $k-1$
(\code{gtzWeighted\_belowWindow\_of\_predecessor}), а меньшие $m$ --- по
предложению~\ref{prop:mono}. Зона пуста при $k=3$, а при $k=4$ это
клетка $(7,4)$.
\item $2k\le m\le k(k+1)/2$ --- \emph{окно}: поднимаемся на один размер за
шаг. Шаг (\code{gtzWeighted\_succ\_of\_hinge\_of\_reach}) требует трёх
вещей: $\GTZ(m-1,k)$; лемма о коллинеарной паре при $(m,k)$; связность
систем без коллинеарных пар с якорем, у которого есть строго подходящая
$k$-ка. Связность доказана при всех рангах $\ge4$ и всех размерах окна
(\code{sharpWindowAnchorReachRankFourAndUp}), при ранге~3 --- икосаэдром.
Доказательство шага --- дословно набросок теоремы~\ref{thm:hinge}.
\item $m>k(k+1)/2$: предложение~\ref{prop:walk}.
\end{enumerate}
Недоказанными остаются только леммы о коллинеарной паре в окне, и ровно
они --- две общие аксиомы реестра:
\begin{itemize}
\item \code{obligationSubThresholdBandHinge}: при $k\ge3$ и $2k\le m<k(k+1)/2$,
из $\GTZ(m-1,k)$ следует, что всякая граничная система размера $m$ ранга $k$
имеет коллинеарную пару;
\item \code{obligationThresholdCellHingeRankFourAndUp}: то же при
$k\ge4$ и $m=k(k+1)/2$.
\end{itemize}
При $k=3$ окно --- одна клетка $m=6$, первая аксиома пуста, а вторая при
$k=3$ --- это и есть лемма из~\S\ref{sub:cell}, уже сведённая к пяти
клеточным аксиомам (\code{obligationThresholdCellHinge}). Отсюда теорема:
\emph{при двух аксиомах об окне утверждение для всех рангов равносильно
одной клетке $\GTZ(6,3)$} (\code{everyRank\_iff\_sixThree}).

Размер окна растёт: $(k-1)(k-2)/2$ клеток --- одна при $k=3$, три при
$k=4$ (размеры $8,9,10$), шесть при $k=5$. Каждая требует своей леммы о
коллинеарной паре.

\subsection{Что переносится с $(6,3)$, а что нет}

\begin{itemize}
\item \emph{Переносится дословно:} теорема~\ref{thm:top} (одна решающая
клетка на ранг, $m=\dim\Sk_k$); монотонность; прогулка по напряжению;
связность; склейка коллинеарной пары; критерий через матрицу Грама и
цепочку миноров (с $k$ условиями вместо трёх); закон присоединённой матрицы
в форме $\adj N=e_{k-1}(N)\,ss\T$ для вырожденной $N$ ранга $k-1$; зеркало;
балансы (с $k$ вместо $3$); воронка.
\item \emph{Переносится с оговорками:} разбор граничных систем на
«с напряжением» и «свободные». В клетке $m=k(k+1)/2$ свободность означает,
что матрицы $g_cg_c\T$ --- базис $\Sk_k$, как при $k=3$. Но в ветви «с
напряжением» при $k=3$ знаки напряжения распадаются ровно $3+3$, и это даёт
разбор; при $k\ge4$ допустимых распадов много, и разбор ветвится.
\item \emph{Не переносится:} в клетках окна ниже $k(k+1)/2$ матрицы
$g_cg_c\T$ не порождают $\Sk_k$, напряжения не обязательны, веса не
определяются направлениями, и инструменты, опиравшиеся на квадратность
системы~\eqref{eq:parseval}, теряют предпосылку. При $k=3$ таких клеток нет;
при $k=4$ это $(8,4)$ и $(9,4)$. Это новая математика, а не повтор.
\item \emph{Обязательно везде:} вещественность. Теорема~\ref{thm:complex}
запрещает бесполевое доказательство при любом $m\ge k+2$.
\end{itemize}

Разумный порядок: закрыть $(6,3)$; затем $(7,4)$ бесплатно (зона~2); затем
$(8,4)$, $(9,4)$, $(10,4)$ --- первые три клетки, где лемма о коллинеарной
паре нужна вне квадратного случая и в нём; затем общая схема по образцу.

% ==========================================================================
\section{О роли поля и о решающем размере}\label{sec:field}

\begin{example}[четыре равноугольные прямые в $\CC^2$]\label{ex:sic}
Пусть $\omega=e^{2\pi i/3}$, $g_0=(\sqrt2,0)$,
$g_j=\bigl(\sqrt{2/3},\ \sqrt{4/3}\,\omega^{j-1}\bigr)$, $j=1,2,3$,
$t_c=1/4$. Тогда $\ell_c=2$ и $|\langle g_a,g_b\rangle|^2=4/3$ при $a\ne b$.
Для любой пары $\tr S_C=4$, $\det S_C=8/3$, и наименьшее собственное значение
каждой из шести пар равно $\alpha_\star=2-2/\sqrt3=0{,}845\ldots<1$.
\end{example}

Положим $\alpha_k(\FF)=\inf\max_{|C|=k}\lambda_{\min}(S_C)$ по системам
над $\FF$ всех размеров; $\GTZ$ при ранге $k$ равносильно $\alpha_k(\FF)\ge1$.
Известно $\alpha_1(\RR)=\alpha_1(\CC)=\alpha_2(\RR)=1$; ромб даёт
$\alpha_3(\RR)\le1$; равенство $\alpha_2(\CC)=\alpha_\star$ доказано
предельным переходом по рациональным весам, в формализации пока только
оценка сверху. Комплексная граничная система ранга один при $(6,3)$
(\S\ref{sec:corner}) показывает, что и при $k=3$ над $\CC$ утверждение
ложно --- в согласии с теоремой~\ref{thm:complex} --- и \emph{локализует}
вещественность: в ветви ранга один она нужна ровно в E2.

Эрмитовы матрицы порядка $k$ над $\FF\in\{\RR,\CC\}$ образуют вещественное
пространство размерности $N_k(\RR)=k(k+1)/2$, $N_k(\CC)=k^2$.
Условие~\eqref{eq:parseval} --- линейная система в этом пространстве с $m$
неизвестными.

\begin{conjecture}\label{conj:law}
Решающий размер при ранге $k$ над $\FF$ есть $m=N_k(\FF)$: при $m<N_k(\FF)$
утверждение следует из теоремы~\ref{thm:corank} и её аналогов, а при
$m=N_k(\FF)$ решается его судьба.
\end{conjecture}

Проверка: $\FF=\RR$, $k=2$: $N_2=3=k+1$, верно. $\FF=\CC$, $k=2$: $N_2=4=k+2$,
ложно с постоянной $\alpha_\star$. $\FF=\RR$, $k=3$: $N_3=6=k+3$ ---
оставшаяся клетка. Над $\RR$ теорема~\ref{thm:top} подтверждает первую
половину гипотезы для всех $k$.

\begin{remark}
Соблазнительно добавить, что при $m=N_k(\FF)$ худшая конфигурация --- $N_k$
равноугольных прямых. Над $\CC$ при $k=2$ это так. Над $\RR$ при $k=3$ шесть
равноугольных прямых --- диагонали икосаэдра, и они подходят \emph{строго};
именно поэтому икосаэдр служит якорем в теореме~\ref{thm:hinge}. Аналогия
работает в одну сторону.
\end{remark}

\begin{thebibliography}{9}

\bibitem{gtz}
С.~А.~Горейнов, Е.~Е.~Тыртышников, Н.~Л.~Замарашкин.
A theory of pseudoskeleton approximations.
\emph{Linear Algebra Appl.} \textbf{261} (1997), 1--21.

\bibitem{airi}
S.~Sengupta, M.~Pautov. arXiv:2604.05944.

\bibitem{main}
Полный текст: \code{paper/main.tex}; формализация: каталоги \code{Gtz/}
и \code{Skeleton/} того же репозитория.

\end{thebibliography}

\vspace{4pt}
\noindent\footnotesize
Утверждения с именами в \code{моноширинном} шрифте формализованы в Lean~4
и проходят проверку ядром без дополнительных аксиом; реестр открытых
аксиом --- файл \code{Skeleton/Obligations.lean}, ровно семь записей.
Утверждения без имени доказаны в тексте. Численные утверждения помечены
словами «численно», «вычисление»; ни одно из них не использовано как
посылка.

\end{document}
